% Dersin Kodu ve Adı (Dönem Yılı ve Mevsimi)
% İMÜ Fakülte ve Bölüm Adı

% İzlence

\documentclass{article}

\newcommand{\course}{Dersin Kodu ve Adı}
\newcommand{\semester}{Dönem Yılı ve Mevsimi}

\usepackage[utf8]{inputenc}
\usepackage[turkish, shorthands = :!]{babel}
\usepackage[T1]{fontenc}
\usepackage{lmodern}

\usepackage{indentfirst} % Indent first paragraph after section header (başlıklardan sonra ilk paragrafın sonrakiler gibi girintilenmesi için)

\usepackage{fancyhdr, tabularx}

\renewcommand{\headrulewidth}{0pt}

\pagestyle{fancy}
\fancyhf{}
\fancyfoot[C]{\thepage} % sayfa alt bilgisinin (footer) orta kısmı
\fancyfoot[R]{\emph{\semester}} % sayfa alt bilgisinin sağ kısmı

\begin{document}

\begin{center}
İstanbul Medeniyet Üniversitesi \\
Fakülte Adı \\
Bölüm Adı \\

\setlength{\fboxsep}{5pt}
\bigskip
\framebox[\linewidth]{\textbf{\course}}
\end{center}

\bigskip

\subsubsection*{Öğretim Üyesi}

\hangindent = \parindent
\hangafter = 1
%
Murat Güngör \\
Göztepe Kuzey Yerleşkesi, F110 \\
murat.gungor@medeniyet.edu.tr \\
0216 280 3333\,/\,4061

%\subsubsection*{Araştırma Görevlisi}
%
%\hangindent = \parindent
%\hangafter = 1
%%
%Seçkin Damar \\
%Göztepe Kuzey Yerleşkesi, H006 \\
%seckin.damar@medeniyet.edu.tr \\
%0216 280 3333\,/\,3443
%Şule Nur Sargın \\
%Göztepe Kuzey Yerleşkesi, H006 \\
%sulenur.sargin@medeniyet.edu.tr \\
%0216 280 3333\,/\,3443

\subsubsection*{Ders Saatleri}

\begin{tabular}{@{}ll@{}} % @ işaretli kısımlar hücre içi sol ve sağ yatay boşlukları ortadan kaldırmak için
Gün, **.**-**.**, B*** \\
\end{tabular}

\subsubsection*{Ders İçeriği}

\hangindent = \parindent
\hangafter = 1
%
İlk konu başlığı, ikinci konu başlığı, \ldots, son konu başlığı.

%\subsubsection*{Kaynaklar}
\subsubsection*{Kaynak}

\hangindent = \parindent
\hangafter = 1
%
İlk yazarın soyadı, ikinci yazarın soyadı, \ldots, son yazarın soyadı. \emph{Kitabın Adı}. Yayınevi, Yıl.

\subsubsection*{Ölçme ve Değerlendirme}

\begin{tabular}{@{}ll@{}} % @ işaretli kısımlar hücre içi sol ve sağ yatay boşlukları ortadan kaldırmak için
Ara Sınav & \%40 \\
Final & \%60 \\
\end{tabular}

\subsubsection*{Yoklama}

\begin{itemize}

\item Öğrencilerin derslere en az yüzde yetmiş oranında devamı zorunludur.
% Daha önce "Öğrencilerin derslere en az yüzde yetmiş oranında devamı zorunludur; başka bir deyişle, devamsızlık hakkı dört buçuk haftadır." yazıyorduk ama dört buçuk hafta kısmını artık hiç karıştırmayalım. Resmî tatillerden dolayı bu azalabilir. Ona göre hesaplayıp hareket etsinler.
% Dönem boyunca ders işlenecek 14 hafta varsa (ara sınav haftası dahil dönem 15 hafta ise) 14 × 0,7 = 9,8 hafta devam zorunluluğu doğuyor. 9½ desek 14 - 9½ = 4½ hafta devamsızlık hakkı oluyor. Dönem boyunca ders işlenecek 13 hafta olsaydı devamsızlık hakkının 4 hafta olduğunu söyleyecektik.

\item Devamsızlık yapan öğrenci, İMÜ Lisans Eğitim-Öğretim ve Sınav Yönetmeliği'nde belirtilen durumlarda mazur sayılır. Bununla birlikte, tam teşekküllü olmayan sağlık kuruluşlarından alınan istirahat raporları da kabul edilecektir.

\end{itemize}

\subsubsection*{Derse Katılım}

\begin{itemize}

\item Tahtaya kalkıp soru çözmek, ders notlarındaki veya slaytlardaki hataları -basit imla hataları dahi olsa- fark edip söylemek ve sınıfta sorulan bazı sorulara cevap vermek suretiyle öğrencilerin derse katılması beklenmektedir.

\item Derse katılan öğrenciler artı alacaktır. Ara sınavdan önce alınan ilk artı 4, ikinci artı 3, üçüncü artı 2, sonraki artılar 1 puan sayılıp ara sınav notuna eklenecektir. Benzer uygulama, ara sınavdan sonraki haftalarda final ve bütünleme için de geçerlidir. Ancak bir sınava eklenecek puan toplamda 20'yi geçemez.

\end{itemize}

\subsubsection*{Sınavlar}

\begin{itemize}

\item Ara sınava bir, finale ve bütünlemeye iki çalışma kâğıdı ile birlikte girilebilir. Çalışma kâğıtları A4 boyutunda olmalıdır. Çalışma kâğıtlarının hem önü hem arkası kullanılabilir.

\item Öğrenciler sınav günleri bilimsel hesap makinelerini yanlarında bulundurmalıdır. Cep telefonu, akıllı saat vb. cihazlar kullanmak yasaktır.

\item Sınavlarda bazı sorular bilgisayar yardımıyla çözülecektir. Dolayısıyla, öğrenciler sınav günleri dizüstü bilgisayarlarını yanlarında getirmelidir. Gerekli programlar bilgisayara önceden yüklenmiş olmalıdır. Bilgisayarlar sınav boyunca mutlaka uçuş modunda tutulacaktır.
% Bilgisayarlı kısımda istatistiksel tabloları toplayalım (eğer klasik kısımda dağıttıysak). Sadece formül kâğıtları ve bilgisayara önceden yükledikleri kaynakları kullanabilsinler.

\item Sınavdan önceki son derste işlenen konular ara sınava dahildir. Finale ve bütünlemeye dönem boyunca işlenen bütün konular dahildir; bununla birlikte, bu sınavlarda genellikle ara sınavdan sonra işlenen konular ağırlıklı olur.

\item Sınavların bilgisayarlı kısmının ağırlığı, bilgisayarın derste kullanım yoğunluğuna göre değişmekle birlikte, \%50'yi geçmemektedir.

\end{itemize}

\bigskip

\noindent
\textit{Not: (1) İletişim için e-posta kullanınız, ÜBYS üzerinden mesaj atmayınız. Gönderdiğiniz e-postalarda adınızı, soyadınızı ve hangi dersin öğrencisi olduğunuzu belirtmeyi unutmayınız. (2) Dönem sonunda notlandırmanın nasıl yapılacağına dair soru sormayınız. Ne gerekçeyle olursa olsun sınava puan ekleme, harf notu aralıklarının değiştirilmesi vb. uygunsuz taleplerde bulunmayınız.}
	
\end{document}